\section{Introduction}\label{sec:intro}
A karambit is a curved, ring-retained utility and defensive knife gripped in a closed fist with the index or little
finger seated through a rear retention ring. Fabricating a one-piece karambit in Ti-6Al-4V by laser powder-bed fusion
(L-PBF) opens an expressive three-dimensional design space: the grip and blade core can be hollowed into a skin over a
triply periodic minimal-surface (TPMS) lattice with integral powder-evacuation ports, while variable-thickness ribs,
finger grooves and dual-edge bevels are built directly into the single monolithic part. Yet unrestricted evolutionary
search in three dimensions also creates severe specification traps. A coarse voxel grid can hide mass or bridge thin
gaps; a beam or plate surrogate can credit interior material that is absent from the printed build; a blade profile can
leave an unground shoulder where a sparse station check does not look; and deep, narrow serration gullets can inflate
apparent edge length while leaving brittle needle teeth that snap under an oblique snag.

Generation~3 (G3) of the Jarambit project addresses these failure modes by formulating the whole knife---retention
ring, ergonomic grip, guard, hollow lattice interior, and a double-edged, double-serrated claw blade---inside a unified
constructive decoder and multi-fidelity finite-element evaluation pipeline. This paper reports the v\gVersion{} release
and the best individual of its cohort of record, \gDesign{} (commit \gCommit, decoder hash \gDecoder).

Three architectural principles govern the v\gVersion{} release:
\begin{enumerate}
\item \emph{Evolution designs the knife from scratch.} Every island in the cohort is initialised exclusively from
scrambled low-discrepancy Sobol sequences (\src{crates/g3/src/g3_genome.rs}{\code{g3_genome.rs}}). No hand-tuned seed,
template CAD model or post-hoc manual edit enters the population. Engineering laws are encoded either as invariant
geometric mappings inside the decoder or as measured physical responses in the objective.
\item \emph{Edge and serration mechanics are intrinsic to the decoder.} Rather than treating sharpening and serrations
as a post-process applied after structural optimisation, decoder revision~r2 enforces five primary-grind invariants
(EDGE-1..5) and four serration-integrity laws (SERR-1..4) directly on the signed distance field
(\src{crates/g3/src/g3_sdf.rs}{\code{g3_sdf.rs}}; Section~\ref{sec:design}). Both cutting edges carry a single primary
bevel at a titanium-appropriate angle (17.5--22.5\textdegree{} per side) measured normal to the true serrated edge
curve, run unbroken to the piercing point, and carry scalloped serrations whose individual tooth roots are checked
against combined in-plane snag (\gToothSnag~N) and cross-blade pry (\gToothPry~N) loads
(\src{crates/g3/src/g3_qd.rs}{\code{serration_tooth_strength}}).
\item \emph{Every reported quantity traces to a single committed measurement receipt.} All scores, dimensions, safety
factors, convergence checks and rankings in this paper and in the repository \code{README.md} are generated
automatically from \code{evidence/g3/v\gVersion/release.json} via
\src{crates/g3/src/bin/g3_release_tex.rs}{\code{g3_release_tex.rs}} and \src{ops/ci/gen_readme.py}{\code{gen_readme.py}}.
Analytic Tier-0 screening estimates never appear in reported results; every quoted value comes from voxel-SDF geometry,
plate and solid finite elements, and an explicit \code{hex8} voxel finite-element certification of the printed part.
\end{enumerate}

The remainder of this paper details the 179-gene representation and edge/serration decoder laws
(Section~\ref{sec:design}), the veto-free multiplicative-ergonomics objective \gMetric{} (Section~\ref{sec:objective}),
the \gIslands-island CVT-MAP-Elites search and three-tier evaluation cascade (Section~\ref{sec:search}), the
multi-grid and voxel finite-element verification chain (Section~\ref{sec:measure}), the evolutionary trajectory and
best individual \gDesign{} (Section~\ref{sec:results}), the explicit physical and modelling limits of the result
(Section~\ref{sec:limits}), and the complete reproduction protocol (Section~\ref{sec:repro}).
